\section{结论}
本文建立了从Geant4沉积总量、TCAD等效源到DRAM读出裕量的可追溯评估链。历史D1结果表明，无粒子读出过程可产生约$12.37\mV$的固定参考偏移；同一信号状态在匹配参考下发生差分符号变化。由此明确，中性参考负责定义读出裕量，同位无粒子对照负责识别粒子增量，存储电荷扰动与最终判据应分别报告。

独立R50版本已完成参考、源隔离、代表点数值检查及32事件的64个训练响应。基于三个沉积与几何特征的四个岭回归器，在八组物理工况留一交叉验证中得到约$0.10$--$0.11\fC$的电荷扰动平均折RMSE及$0.96$--$1.01\mV$的裕量损失平均折RMSE，均低于对应训练均值基准。该结果支持当前条件响应代理的训练内可行性；独立插值与尺寸外推仍需按冻结预测完成验证。

AI的已证实作用是辅助构建受约束输入、提取与审计流程，未证明任意自然语言自主仿真或操作效率优势。研究结论限于未校准1T1C探索模型、固定$5\,\mathrm{MeV}$粒子条件和当前电磁输运及二维近似，不推导真实器件误码概率、核反冲效应或商业节点可靠性。后续工作应以完成既定独立测试、边界检查和参数敏感性归档为依据扩充结论，避免用计划或训练拟合代替验证。
